\chapter{Randgitter, Dynamik und lokale Rekonstruktion}

\section{Ausgangspunkt und Geltungsbereich}

Die späteren Arbeiten zur Rand-, Gitter- und Eichstruktur bilden eine logisch zusammenhängende Kette. Ihr Ausgangspunkt ist \textbf{kein} aus den TFPT-Prinzipien bereits abgeleiteter Materiesektor, sondern ein präzise deklarierter Kandidat: zehn chirale bosonisierte Kanäle mit

\[
\Gamma=\mathbb Z^{9,1},\qquad K=\operatorname{diag}(1^9,-1),
\]

also neun rechts- und einem linksbewegten Kanal. Die Resultate entscheiden innerhalb dieses Kandidaten sehr scharf, welche Nullvektoren, Ladungen, lokalen Felder und Erweiterungen miteinander verträglich sind. Sie entscheiden ebenso scharf, welche scheinbar naheliegenden Identifikationen scheitern. Sie leiten aber weder diese zehn Kanäle noch ihre Energie, ihren Zustand oder ihre Eichdynamik aus P1/P2 her. Diese Grenze wird in allen folgenden Aussagen beibehalten.

Im Folgenden bedeutet „exakt“ eine algebraisch bewiesene Aussage unter diesen festgehaltenen Voraussetzungen; „bedingt“ markiert eine zusätzlich gewählte Energie, Eichverbindung oder Sektorregel; „offen“ bezeichnet den noch fehlenden Transport aus derselben ursprünglichen Quelle.

Der zentrale Nullvektor ist

\[
n=(1,1,1,-1,-1,-1,-1,-1,-1,3),
\]

mit \(n^TKn=0\). Daneben steht der gewöhnliche bilineare Nullvektor

\[
z=e_9-e_{10}.
\]

Für die ersten acht Koordinaten wird

\[
a=(1,1,1,-1,-1,-1,-1,-1)
\]

verwendet. Zwei Einbettungen spielen später verschiedene Rollen:

\[
T(r)=\bigl(r,-a\!\cdot\!r/2,a\!\cdot\!r/2\bigr),
\]

\[
F(r)=T(r)-\frac{a\!\cdot\!r}{2}n
=\bigl(r-(a\!\cdot\!r)a/2,0,-a\!\cdot\!r\bigr).
\]

Diese beiden Formeln sind algebraisch eng verwandt, aber auf geladenen Feldern nicht austauschbar. Genau diese Unterscheidung wird zum roten Faden der späteren Ergebnisse.

\section{Auswahl des Randkandidaten und erste physikalische Markierungen}

\src{37} beginnt mit allen primitiven charakteristischen Nullvektoren minimaler euklidischer Länge im Gitter \(I_{9,1}\). Die Charakteristik erzwingt ungerade Komponenten; Minimalität gibt \(|n_{10}|=3\). Die globale Quellenladung, der Glue und die orientierte Farb-/Schwach-/Familienzerlegung reduzieren die Kandidaten stufenweise. \src{38} ersetzt dabei eine zunächst separat angenommene Blockpermutationssymmetrie durch die tatsächliche Gruppenbedingung

\[
G_\mathrm{SM}=\frac{SU(3)\times SU(2)\times U(1)}{\mathbb Z_6}=S(U(3)\times U(2)).
\]

Für einen farb- und schwachsingulären Zusatzkanal mit Hyperladung \(y\) lautet die Abstiegsbedingung \(y\in\mathbb Z\). Zusammen mit der Quellmarkierung wählt sie eindeutig den oben angegebenen orientierten Vektor \(n\) und \(y=1\).

Der zusätzliche Ladungstest in \src{37} macht diese Bedingung unmittelbar sichtbar. Mit

\[
Y_8=(-1/3,-1/3,-1/3,1/2,1/2,0,0,0),\qquad
Y_{10}=(Y_8,y,y)
\]

gilt

\[
Y_{10}\cdot n=-2+2y.
\]

Die Wechselwirkung ist daher in der festgelegten vektorartigen Singulettklasse nur für \(y=1\) hyperladungsneutral. Dann bleibt sogar das ganze transportierte E8-Ladungswörterbuch erhalten. In der orthogonalen Zerlegung

\[
\Gamma=F(E_8)\oplus \mathbb Ze_R\oplus\mathbb Zm,
\qquad e_R=-e_9,\qquad m=n+e_9,
\]

tragen \(e_R\) und \(m\) Hyperladungen \(-1\) und \(+1\); ihr neutraler Summenvektor ist \(n\). Dies ist ein exakter Ladungstest für das massive Komplement, keine Herleitung eines neuen vierdimensionalen Elektrons.

Der native E8-Kopplungstensor besitzt damit zwar die richtigen geladenen Down-/Lepton-Slots, löst aber das Flavorproblem nicht automatisch. Aktiviert man in einem festen neutralen Kanal nur ein Familientriplet \(h\), entsteht

\[
A(h)_{ij}=\epsilon_{ijk}h_k,\qquad A(h)h=0,
\qquad \operatorname{rank}A(h)=2.
\]

Ein zusätzlicher Singulettkanal \(k\) kann den erweiterten Block

\[
Y(h,k)=\begin{pmatrix}A(h)&k\\-k^T&0\end{pmatrix},
\qquad \det Y=(h^Tk)^2
\]

algebraisch auf vollen Rang vier bringen. Dieser bedingte Weg wählt weder \(k\) noch einen schweren vierten Zustand, drei leichte Familien oder ihre Hierarchie aus. Er ist eine genaue Anforderung an die fehlende Quellkopplung, kein fertiger Yukawa-Mechanismus. \src{38}

Am zusätzlich gewählten Punkt

\[
V_\mathrm{aux}=K+2Kmm^TK
\]

zerfällt die Energie in einen E8-Sektor und ein Dirac-Paar. Erst zusammen mit dem zusätzlich gewählten nichtverschwindenden Cosinus-Term zum Nullvektor $n$ wird das Paar am bezeichneten Rekonstruktionspunkt massiv. Ohne diesen Term ist die Energiematrix allein keine Massenerzeugung. Dort hat \(n\) Skalendimension eins, während \(z\) Dimension neun besitzt. \src{40} zeigt jedoch, dass dieser Punkt eine echte Änderung der Energie ist. Der passive Basiswechsel der ursprünglichen freien Metrik \(V_0=I\) liefert \(W^TV_0W\), nicht automatisch \(V_\mathrm{aux}\). Besonders deutlich ist die Verteilung der 240 E8-Wurzelvertices: Am freien Punkt besitzen 56, 112, 56 und 16 von ihnen die Dimensionen 1, 2, 5 und 10; erst bei \(V_\mathrm{aux}\) haben alle Dimension eins. Die E8-Energie ist deshalb keine bloße Umbenennung der freien QWZ-Quelle.

Auf einem ausdrücklich deklarierten Interpolationspfad zwischen beiden Energien sind \(n\) und \(z\) die einzigen neutralen selbstnullen Vertices, die relevant oder marginal werden können. Alle anderen zulässigen neutralen Nulloperatoren bleiben auf dem gesamten Pfad bei

\[
\Delta\geq \frac72.
\]

Der Mittelpunkt hat \(\Delta(n)=\Delta(z)=2\). Dieser starke Auswahlfilter ist exakt für den angegebenen Pfad. Er ist kein RG-Fluss und kein Beweis, dass die Quelle diesen Pfad oder einen kritischen Punkt auswählt.

\section{Das vollständige Gitter trägt die Statistik}

\src{39} liefert die entscheidende lokale Struktur. Der gemeinsame orthogonale Kern von \(n\) und \(z\) ist

\[
L_0=T(D_8),
\]

während die neutrale Ebene \(P_0=\mathbb Zn+\mathbb Zz\) die Gram-Matrix

\[
\begin{pmatrix}0&2\\2&0\end{pmatrix}
\]

besitzt. Das Untergitter \(L_0\oplus P_0\) hat Index vier in \(\Gamma\). Zwei Glue-Erzeuger können als

\[
f=e_1=T(e_1)+z/2,\qquad
b=F(s)=T(s)+n/2,\qquad s=(1/2)^8,
\]

gewählt werden. Die vier Klassen \(0,f,b,f+b\) entsprechen den \(D_8\)-Diskriminantenklassen \(0,v,s,c\). Der E8-Quotient und der freie \(I_8\)-Quotient verwenden denselben \(D_8\)-Kern, aber verschiedene Erweiterungsklassen.

Die scheinbar „wegprojizierte“ neutrale Ebene ist dabei nicht bedeutungslos. In einer typischen Paarung liefern der \(D_8\)-Anteil und der komplementäre Anteil jeweils eine Halbphase; beide Minuszeichen multiplizieren sich im vollständigen Quellgitter zu lokaler Monodromie \(+1\). Das massive Komplement stellt also die fehlende Lokalitätsphase wieder her. Entfernt man es, darf man die verbleibenden E8-Gewichte nicht automatisch als dieselben lokalen Felder lesen.

Die mikroskopische Parität ist besonders einfach. Für

\[
x=F(r)+Ae_R+Bm
\]

gilt

\[
(-1)^F=(-1)^{x^TKx}=(-1)^{A+B}.
\]

Reine E8-Vertices sind gerade. Jedes ungerade lokale Feld trägt einen ungeraden Anteil des massiven Paares. Die ursprünglichen acht Fermionen verschwinden dabei nicht; sie faktorisieren exakt. Für geeignete Spinorwurzeln \(r_i\) gilt

\[
e_i=F(r_i)-a_i m,\qquad i\leq 8.
\]

Ein ursprüngliches ungerades Fermion ist somit ein E8-Spinorstrom zusammen mit einem massiven ungeraden Vertex. Unterhalb der schweren Lücke annihiliert die niedrige Projektion alle ungeraden lokalen Felder. Eine gerade Schur-Elimination kann keinen fehlenden ungeraden Hilbertraumsektor erzeugen. Das bedeutet nicht, dass die vollständige Theorie keine Fermionen besitzt; es bedeutet nur, dass im faktorisierten Niedrigenergiesektor keine genuin leichten ungeraden lokalen Felder verbleiben.

Diese Aussage korrigiert zugleich einen früher erwogenen Flavorweg. Ein Rang-ergänzender Schur-Block kann algebraisch eine singuläre Dreifamilienmatrix vervollständigen. Im hier ausgewählten faktorisierten Sektor wären direkte Mischungen zwischen einem reinen leichten E8-Feld und genau einem schweren ungeraden Fermion jedoch paritätsverboten. Der abstrakte Rangmechanismus bleibt richtig, ist aber kein vorhandener lokaler Yukawa-Term dieser Phase.

Auch die Determinantenphase muss vollständig behandelt werden. Für endliche invertible Blöcke gilt

\[
\det D=\det D_{QQ}\,
\det\!\left(D_{PP}-D_{PQ}D_{QQ}^{-1}D_{QP}\right).
\]

Am exakt faktorisierten Punkt ist \(Z[j;\lambda]=Z_{E_8}[j]Z_m[\lambda]\). Der massive Faktor gehört zur vollen Wirkung, kürzt sich aber aus einer relativen Down-/Lepton-Antwort

\[
\Xi(\theta)=\frac{Z_d(\theta)Z_e(0)}{Z_e(\theta)Z_d(0)}
\]

exakt heraus, wenn beide Kanäle denselben \(Z_m(\theta)\) tragen. Außerdem ist die Determinante eines freien konstantphasigen Dirac-Kerns

\[
h_\theta(k)=k\sigma_z+m\cos\theta\,\sigma_x+m\sin\theta\,\sigma_y
\]

wegen unitärer Äquivalenz unabhängig von \(\theta\). Eine relative Flavorphase benötigt daher kanalverschiedene Massendeformationen, eine echte E8-Komplementkopplung oder zusätzliche Hintergrund-/Regulatordaten. Eine gemeinsame freie Dirac-Massenphase genügt nicht.

\section{Lokale Rekonstruktion am bedingten Konkurrenzpunkt}

\src{41} untersucht den separat gewählten Punkt

\[
u=(n+z)/2,\qquad v=(n-z)/2,\qquad
V_c=K+2Kvv^TK.
\]

Hier gilt \(\Delta(n)=\Delta(z)=1\). Dieser Punkt liegt nicht auf dem direkten Energiepfad des vorigen Abschnitts; schon der Matrixeintrag \((V_c)_{9,10}=4\), gegenüber null auf dem ganzen direkten Pfad, zeigt die zusätzliche kollektive Dichterichtung.

Die vier Glueklassen besitzen am Punkt \(V_c\) die exakten minimalen Dimensionen

\[
\begin{array}{c|cccc}
\text{Klasse}&0&v&s&c\\ \hline
(-1)^F&+&-&+&-\\
\Delta_{\min}&0&3/4&5/4&3/2.
\end{array}
\]

Unter der zu \(D_5+A_3\) gehörenden \(\mathrm{Spin}(10)\times SU(4)\)-Wirkung verzweigt die gekreuzte Klasse als

\[
c=(16,\bar4)+(\bar{16},4).
\]

Für jedes passende Halbgewicht \(p\) der \(c\)-Klasse sind

\[
x_R=T(p)+u,\qquad x_L=T(p)+v
\]

ganzzahlige, mikroskopisch ungerade Felder. Explizite Vertreter besitzen \(q=3\), \(Y=1/3\) und \(\Delta=3/2\). Der neutrale Anteil ist unverzichtbar: \(u\) und \(v\) allein sind keine lokalen Vektoren des ursprünglichen Gitters. Die Bezeichnungen rechts/links betreffen die übernommene Zerlegung; insbesondere kann ein Gesamtfeld Gewichte \((1,1/2)\) tragen und ist dann nicht rein chiral. Dies ist eine 1+1-dimensionale Feldkonstruktion, keine Identifikation mit einem vierdimensionalen Weyl-Fermion.

Eine konkrete Korrektur ist in der Endfassung festgehalten: Der rohe Vertreter \(f+b\) hat \(\Delta=5/2\), nicht \(3/2\). Die kleinere Dimension wird erst durch Verschiebung um einen echten \(T(D_8)\)-Gittervektor innerhalb derselben Klasse erreicht.

Die bedingte Refermionisierung der beiden Cosinus-Terme führt formal zu zwei Majoranamassen proportional zu \(g_n\pm g_z\). Die zugehörigen halbintegralen Fermionvektoren sind jedoch keine unabhängigen lokalen Felder des UV-Gitters. \src{41} konstruiert einen sechssektorigen, monodromielokalen Kandidaten in \((D_8)_1\times\mathrm{Ising}\), einschließlich des fermionischen einfachen Stroms \(J=(c,\psi_R)\). Exakt sind die UV-Glueklassen, ihre Gewichte und die Monodromietabelle. Nicht hergeleitet sind Massenvorzeichen, Spinstruktur, Defektlinien und der physische IR-Projektor. \src{41} zeigt zusätzlich, dass keine nichttriviale Potenz der vorhandenen festen \(C\)- oder \(J\)-Clock zugleich \(V_c\), den \(z\)-Term und die Hyperladung erhält. Der kritische Kandidat ist daher ein präziser Zielanschluss, keine von der Quelle ausgewählte Phase.

\finding{Lokale Statistik liegt im vollständigen Glue}{Die E8-Gewichte allein tragen nicht die ganze lokale Feldinformation. Das neutrale Komplement stellt die gegenseitige Lokalität wieder her und bindet jedes ungerade Feld an einen geladenen oder massiven Sektor.}

\chapter{Gruppenobstruktion, Eichladung und geladene Lifts}
\section{Die zentrale Gruppenobstruktion und die nichtlineare Randlösung}

Die Quellprüfung \src{44} untersucht zuerst den naheliegenden Ursprung im nativen 64-CAR/60-CCR-Fockmodell. Dort ist die Nicht-Gauß-Struktur wirklich Teil des Zeitgenerators,

\[
H_W=\Delta N_b+g\sum_{A=1}^{60}
\left(b_A^\dagger P_A+P_A^\dagger b_A\right),
\qquad P_A=\sum_{i<j}W_{A,ij}f_jf_i,
\]

mit \(WW^T=8I_{60}\). Dies macht den nativen Fockpfad zu einem ernsthaften Ursprungskandidaten, liefert aber noch keinen Raumträger und keinen Transfer auf \(n,z,V_c\). Dessen Fermionen tragen \((\bar{16},\bar4)\), seine 60 Bosonkanäle \((10,6)\). Ein einziges diagonales Zentrumselement \(\zeta\in\mathrm{Spin}(10)\times SU(4)\) wirkt auf allen nativen Erzeugern mit \(+1\), auf den benötigten ungeraden \(c\)-Feldern aber mit \(-1\). Daher wirkt \(\zeta\) auf der gesamten nativen Operatoralgebra trivial. Ein äquivarianter Transfer eines \(c\)-Felds müsste gleichzeitig \(O=+O\) und \(O=-O\) erfüllen und ist folglich null. Der Ausschluss gilt für beliebig komplizierte native Operatorprodukte, nicht nur bis zu einem getesteten Grad.

Die vor der Familienkompression vorhandenen Clifford-Matrizen besitzen tatsächlich die kovariante Struktur

\[
6\otimes\bar4\longrightarrow4.
\]

In der exportierten geraden Familienkompression verschwinden die sechs einzelnen Übergänge jedoch:

\[
P_\mathrm{even}\gamma_aP_\mathrm{even}=0.
\]

Hilfszustände des Clifford-Aufbaus sind also noch keine ausgeführten physikalischen Feldoperatoren. \src{45} verschärft dies: Fügt man die sechs Vektor-Majoranas als gewöhnliche CAR-Faktoren hinzu, verknüpfen zwei zentrale Zeichen Gruppenwirkung und Gesamtparität so, dass der gewünschte Spinorsektor weiterhin die falsche Statistik besitzt. Auch die volle endliche Clifford-Algebra \(\mathrm{End}(S)\) enthält nur ganzzahlige Differenzgewichte, keine Spinoroperatoren mit halbzahligen Gewichten.

Der ebenfalls geprüfte rationale Austausch \(n\leftrightarrow z\) löst diese Obstruktion nicht. Er ist nur auf einem gemeinsamen Untergitter vom Index zwei integral; dieses enthält die Klassen \(0,c\), aber nicht \(v,s\). Ein formaler Labeloperator mit \(D^2=1+\eta\) ist damit algebraisch formulierbar. Eine lokale Defektalgebra, eine Sektorenregel und ein zeitlicher Intertwiner folgen daraus nicht. Insbesondere ist diese Index-zwei-Struktur nicht mit dem vorhandenen metaplektischen Half-Deck gleichzusetzen. \src{44}

Der gleiche Contract liefert aber eine positive, nichtlineare Randrekonstruktion. Im bedingten \(\Gamma\)-Rand ist

\[
e_9=T(-a/2)+v
\]

ein vorhandenes ursprüngliches Fermion und zugleich ein Keim der \(c\)-Klasse. Die Nullmoden der zusammengesetzten \(T(D_8)\)-Ströme erzeugen daraus den vollständigen 128-Gewichte-Orbit

\[
x_p=T(p)+v.
\]

Die kleinere Produktgruppe erreicht jeweils 64 Gewichte der Zweige \((16,\bar4)\) und \((\bar{16},4)\); erst zusätzliche \(D_8\)-Wurzeln verbinden beide. Entscheidend ist, dass 56 der 112 \(D_8\)-Wurzelströme in den ursprünglichen Fermionkoordinaten quartisch sind. Die Gruppenwirkung ist deshalb keine lineare Bogoliubov-Wirkung der endlichen Quell-CAR-Algebra. Am freien Punkt besitzen die Orbitfelder unterschiedliche Dimensionen; erst die zusätzlich gewählte Dynamik \(V_c\) macht sie zu einem gemeinsamen \(\Delta=3/2\)-Multiplett. Damit ist die lokale Rekonstruktion exakt, die Herkunft ihrer gemeinsamen Dynamik aber weiterhin offen.

\section{Gaußreduktion, kompatible Ladung und Flussindex}

\src{42} analysiert den vorhandenen kompakten U(1)-Rotorparenten, ohne ihn mit der zehnkanaligen Quelle gleichzusetzen. Auf einem endlichen verbundenen Graphen lautet die Gaußbedingung

\[
BE=-\rho,\qquad \rho_x=N_x-1.
\]

Für einen gewählten Spannbaum ist jede Lösung eindeutig

\[
E=E_T(\rho)+Cw,\qquad w\in\mathbb Z^{b_1}.
\]

Die exakte elektrische Energie zerfällt in

\[
|E|^2=\rho^T(BB^T)^+\rho+(w+\theta(\rho))^TQ(w+\theta(\rho)).
\]

Dies bewahrt den longitudinalen Dichtekern und den affinen ganzzahligen Windungsrest. Der zugehörige unitäre Gaußabbau transportiert zugleich sämtliche ursprünglichen Hoppings, Zustände und dieselbe Zeitentwicklung. Das ist ein positiver, vollständiger Reduktionssatz für diesen Rotorparenten. Es ist noch keine Abbildung auf die QWZ-/Clock-Quelle oder auf den E8-Rand.

Jedes endlich getragene Gauß-invariante Observable ist in der Materie-Fermionparität gerade. Ein einzelnes geladenes Fermion ist kein lokales Gauß-invariantes Observable; Wilsonlinien liefern bilineare Transporte oder geladene Sektorintertwiner. Deshalb kann die ungraduierte Observable-Algebra nicht unmittelbar die volle lokale CAR-Feldalgebra aufnehmen. Ebenso scheitert die direkte Identifikation der alten globalen Ladung \(q=\sum_i x_i\) als Eichladung: Unter den E8-Wurzeln treten die Ladungen

\[
q=-4,-2,0,2,4
\]

mit Multiplizitäten \(1,56,126,56,1\) auf. Nur \(E_7\oplus u(1)\) zentralisiert diese Ladung; 114 E8-Schrittoperatoren wären geladen.

Konstruktiv lässt sich dagegen jede ganzzahlige Ladung klassifizieren, die den gesamten \(F(E_8)\)-Sektor neutral lässt. Sie hat die Form

\[
g(v,u)=(va,u,-3v).
\]

Die Erhaltung bereits vorhandener \(q\)- oder \(Y\)-Markierungen erzwingt \(u=-v\). Primitive Normierung liefert eindeutig bis auf Vorzeichen

\[
\boxed{g=Kn=(a,-1,-3).}
\]

Diese Ladung ist rein und gemischt anomaliefrei in der bedingten 1+1-dimensionalen Bilinearrechnung:

\[
g^TK^{-1}g=q^TK^{-1}g=Y^TK^{-1}g=0.
\]

Die gravitative Signatur bleibt jedoch acht. Zudem gilt \(g\cdot n=0\), aber \(g\cdot z=2\). Eine wirkliche Eichung erlaubt daher den nackten \(n\)-Term, nicht zugleich den nackten konkurrierenden \(z\)-Term. Die lokalen \(c\)-Spinorfelder tragen Ladung \(\pm1\), während alle E8-Ströme neutral sind. Der Nachtrag \src{42} zeigt: Das Hinzufügen eines chemischen Terms \(\lambda Q_g\) kann die bereits vorhandenen unterschiedlichen Clock-/Quellfrequenzen der E8-Wurzeln nicht ausgleichen, weil \(Q_g\) auf ihnen verschwindet.

\src{43} setzt nun ausdrücklich eine zusätzliche physikalische Prämisse: eine kompakte U(1)-Verbindung, die an genau diese zehn chiralen Kanäle mit Ladung \(g\) gekoppelt ist. Für Einheitsfluss ist der signierte Index

\[
I=Kg=n
\]

bis auf die gemeinsame Orientierungswahl. Der minimale lokale Sättigungsoperator enthält zwölf Fermionfaktoren. Die drei identischen Faktoren im Ladung-drei-Kanal benötigen die Ableitungsordnungen 0, 1 und 2; deshalb hat der freie Operator \((h,\bar h)=(9/2,9/2)\) und \(\Delta=9\). Der Index bestimmt die Auswahlregel eines möglichen 't-Hooft-Vertices. Er bestimmt weder dessen Koeffizienten noch eine lokale Cosinus-Exponentiation, einen Vakuumwinkel, eine Phase oder den Fluss-Summenmaßstab.

Auf Gitterebene ist dagegen exakt

\[
\ker(g\cdot)=F(E_8)\oplus\mathbb Zn,\qquad
\ker(g\cdot)/\mathbb Zn\cong E_8.
\]

Der Quotient ist sogar vollständig als \(A_8\)-Index-drei-Erweiterung darstellbar: 72 bilineare und je 84 vierfaktorige positive beziehungsweise negative Wurzelvertreter ergeben die 240 Norm-zwei-Klassen. Ihre Quotientennorm macht sie am ursprünglichen freien Punkt noch nicht zu chiralen Gewicht-eins-Strömen.

\section{Geladene Lifts und der minimale lokale Reparaturraum}

Der neutrale Quotient legt die Wirkung auf geladene Felder nicht fest. \src{46} klassifiziert deshalb alle ganzzahligen additiven isometrischen Schnitte desselben Quotienten:

\[
S_\lambda(r)=F(r)+(\lambda\cdot r)n,\qquad \lambda\in E_8.
\]

Die gewünschte gemeinsame Wirkung auf den geladenen Keim \(e_9\) würde dagegen

\[
\lambda=a/2
\]

erzwingen. Dieser Vektor liegt im gegenüberliegenden Spinor-Coset und nicht im gewählten E8-Gitter. Folglich existiert kein integraler voller E8-Schnitt mit derselben exakten Cartanpaarung auf \(e_9\). Auf \(D_8\) funktioniert \(T\) vollständig; genau dort ist die maximale fortsetzbare Unterstruktur. Die 112 D8-Wurzeln sind integral, alle 128 E8-Halbwurzeln verlangen halbe Vielfache von \(n\).

Die stärkere Lokalitätsformel lautet für jede Halbwurzel \(r\) und jedes ursprüngliche Feld \(x\)

\[
e^{2\pi iB(T(r),x)}=(-1)^{g\cdot x}.
\]

Auf jedem ungeraden Originalfeld entsteht also nichttriviale Monodromie. Ein Kokzyklus kann Austauschzeichen organisieren, diese volle Monodromie aber nicht entfernen. Die globalen Markierungen \(q,Y\) wählen den Lift nicht aus, weil sie \(n\) annihilieren; verschiedene \(S_\lambda\) sind durch integrale Eichler-Transvektionen erreichbar. Werden Gitter, Felder, Energie und Zustand gemeinsam transformiert, sind sie Koordinatenwechsel. Hält man den Quellkeim oder die Zeit fest, sind sie verschiedene physische Wörterbücher.

Der minimale lokale Reparaturraum für den vollen \(T(E_8)\)-Lift ist ein anderer Nachbar des ursprünglichen ungeraden Gitters. Setze

\[
\Gamma_0=\{x\in\Gamma:g\cdot x\in2\mathbb Z\},\qquad h=n/2,
\]

\[
L_h=\Gamma_0+\mathbb Zh.
\]

Dann ist \(L_h\) gerade, unimodular und

\[
L_h=T(E_8)\oplus\mathbb Zh\oplus\mathbb Zz
\cong E_8\oplus U=II_{9,1}.
\]

Dies ist die kleinste ganzzahlige Erweiterung, die den vollen \(T(E_8)\)-Lift bei Erhalt des gesamten geraden Originalsektors enthält. Sie schließt aber alle ursprünglichen ungeraden Felder aus: \(L_h\cap\Gamma=\Gamma_0\), insbesondere \(e_9\notin L_h\). Der Diskriminantengruppe zufolge gibt es zwei gerade Nachbarn, \(L_h\) und \(L_c=\Gamma_0+\mathbb Z(h+e_9)\); erst die gewünschte \(T(E_8)\)-Einbettung wählt bedingt \(L_h\). Auf diesem neuen Gitter ist \(g/2\), nicht \(g\), die primitive Ladung. Aus \(h=n/2\) folgt daher kein physisches Halbflussereignis der alten Theorie.

